\par\Needspace{8\baselineskip}\bigskip
\SourceLeaf{457}
\pdfbookmark[3]{Darstellung, 9}{pallas-II-article9}
\begin{center}9.\end{center}
\noindent\hspace*{\parindent}Nehmen wir an, dass die Störungen die Größen \(\xi,\eta,\zeta\) zu den Grundgleichungen hinzufügen, sodass wir erhalten
\[
\begin{aligned}
0&=\dfrac{\mathrm{dd}x}{\mathrm{d}t^2}+\dfrac{mx}{r^3}+\xi\\[5pt]
0&=\dfrac{\mathrm{dd}y}{\mathrm{d}t^2}+\dfrac{my}{r^3}+\eta\\[5pt]
0&=\dfrac{\mathrm{dd}z}{\mathrm{d}t^2}+\dfrac{mz}{r^3}+\zeta.
\end{aligned}
\]
\par Man sieht, dass \(\xi,\eta,\zeta\) die Kräfte sind, die sich ergeben, wenn man alle Kräfte, die die relative Bewegung des Planeten gegenüber der Sonne stören, nach den Richtungen zerlegt, die denen der drei Achsen entgegengesetzt sind. Man kann sich diese Kräfte auch zu einer einzigen Kraft vereinigt denken, deren Betrag \(=\chi\) ist und deren Richtung parallel zu einem Radius unserer Kugel verläuft, der vom Mittelpunkt zum Punkt \(Q\) auf der Oberfläche gezogen ist. Dann gilt
\[
\begin{aligned}
\xi&=-\chi\cos XQ\\
\eta&=-\chi\cos YQ\\
\zeta&=-\chi\cos ZQ.
\end{aligned}
\]
\par Es kann nützlich sein, statt der Kräfte \(\xi,\eta,\zeta\) diejenigen einzuführen, die sich bei der Zerlegung von \(\chi\) nach drei anderen, zueinander senkrechten Richtungen ergeben. Nehmen wir zunächst an, dass die erste senkrecht zur Ebene der oskulierenden Bahn wirkt, die zweite in Richtung des Radiusvektors und folglich die dritte in der Bahnebene senkrecht zum Radiusvektor. Genauer gesagt, nehmen wir an, dass diese drei Komponenten parallel zu den drei Kugelradien wirken, die zu den Punkten \(M,L,\Lambda\) gezogen sind. Damit ergeben sich
\[
\begin{aligned}
\chi\cos MQ&=-\xi\cos XM-\eta\cos YM-\zeta\cos ZM\\
\chi\cos LQ&=-\xi\cos XL-\eta\cos YL-\zeta\cos ZL\\
\chi\cos\Lambda Q&=-\xi\cos X\Lambda-\eta\cos Y\Lambda-\zeta\cos Z\Lambda.
\end{aligned}
\]
\par Die erste ist also
\[
=-\dfrac{A\xi+B\eta+C\zeta}{\BareRoot{mp}}\mathpunct{\raisebox{.45ex}{,}}
\]
\SourceLeaf{458}die zweite
\[
=-\dfrac{x\xi+y\eta+z\zeta}{r}\mathpunct{\raisebox{.45ex}{,}}
\]
die dritte
\[
=-\dfrac{(Bz-Cy)\xi+(Cx-Az)\eta+(Ay-Bx)\zeta}{r\BareRoot{mp}}\mathpunct{\cdot}
\]
Wir bezeichnen diese drei Kräfte mit \(mW,mT,mV\).
